% Generated draft: requires mathematical and source review.
\documentclass[11pt]{article}
\usepackage[margin=25mm]{geometry}
\usepackage{fontspec}
\setmainfont{Latin Modern Roman}
\usepackage{amsmath,amssymb,mathtools,unicode-math}
\setmathfont{Latin Modern Math}
\usepackage{enumitem,needspace,xurl,hyperref}
\setlength{\parindent}{0pt}
\setlength{\parskip}{5pt}
\setlength{\emergencystretch}{3em}
\newcommand{\N}{\mathbb{N}}
\newcommand{\Z}{\mathbb{Z}}
\newcommand{\Q}{\mathbb{Q}}
\newcommand{\R}{\mathbb{R}}
\newcommand{\C}{\mathbb{C}}
\newcommand{\F}{\mathbb{F}}
\newcommand{\A}{\mathbb{A}}
\newcommand{\PP}{\mathbb{P}}
\newcommand{\E}{\mathbb{E}}
\newcommand{\eps}{\varepsilon}
\newcommand{\dd}{\,\mathrm{d}}
\begin{document}
\section*{UnsolvedMath ID: 30001814}
\textbf{Source problem code:} catalogue entry
\subsection{Definitions and mathematical statement}
Let \(\Gamma\) be a finitely generated infinite group. A partial order on \(\Gamma\) is a relation \(\preceq\) that is reflexive, antisymmetric, and transitive: for all \(x,y,z\in\Gamma\),
\[
 x\preceq x,\qquad
 (x\preceq y\ \text{and}\ y\preceq x)\Longrightarrow x=y,\qquad
 (x\preceq y\ \text{and}\ y\preceq z)\Longrightarrow x\preceq z.
\]
It is a total order if, additionally, every \(x,y\in\Gamma\) satisfy \(x\preceq y\) or \(y\preceq x\). A partial or total order is left-invariant if
\[
 x\preceq y\quad\Longleftrightarrow\quad gx\preceq gy
 \qquad(g,x,y\in\Gamma),
\]
and bi-invariant if
\[
 x\preceq y\quad\Longleftrightarrow\quad gxh\preceq gyh
 \qquad(g,h,x,y\in\Gamma).
\]
\textbf{Question as supplied.} For a finitely generated infinite group \(\Gamma\), does there exist an invariant order relation on \(\Gamma\)? Consider separately the four combinations
\[
\begin{array}{c|cc}
 & \text{partial order} & \text{total order}\\ \hline
\text{left-invariant} & \text{existence?} & \text{existence?}\\
\text{bi-invariant} & \text{existence?} & \text{existence?}
\end{array}
\]
\textbf{Unresolved qualifications.} Under the definitions above, the equality relation is an admissible bi-invariant partial order. The supplied record does not specify whether partial orders must be nontrivial, for example by requiring distinct comparable elements. Such a requirement is therefore not imposed here, but its absence prevents identifying a substantive intended partial-order question. The record also does not specify whether the target is universal existence for all finitely generated infinite groups, an existence criterion for a given such group, or a question restricted to a particular class of groups suggested by the workshop context. No additional restriction is inferred from that context.

The supplied literature assessment says that the report distinguished partial and bi-invariant orders and stated that nonexistence of a bi-invariant partial order had not been proved even for \(\mathrm{SL}(3,\Z)\), the group of determinant-one \(3\times3\) integer matrices. This is retained as a supplied source qualification, not as an independently verified assertion. In particular, its intended meaning requires clarification of the partial-order convention.
\subsection{Short English statement}
The question asks about left-invariant and bi-invariant orders on a finitely generated infinite group, considering both total and partial orders. The intended group quantifier and any nontriviality requirement for partial orders remain unclear. The supplied open-status assessment requires independent review.
\subsection{Sources}
\begin{itemize}
\item Local UnsolvedMath registry, supplied complete source record: permanent ID \(30001814\), title \emph{Invariant Orders on Finitely Generated Infinite Groups}, source problem code \texttt{catalogue entry}, in the \emph{Oberwolfach Reports Open Problems} collection. The record gives proposed year 2011 and names no individual proposer.
\item Supplied report citation: \emph{Arithmetic Groups vs. Mapping Class Groups: Similarities, Analogies and Differences} (2011), Oberwolfach Reports, pp.~1637--1708, DOI \texttt{10.4171/owr/2011/30}. Supplied source URL: \url{https://doi.org/10.4171/owr/2011/30}. The background describes these reports as documenting mathematical workshops and their open-problem discussions.
\item Supplied literature-source entry, labeled ``original Oberwolfach report'': \url{https://doi.org/10.4171/owr/2011/30}. Literature checked at: \texttt{2026-08-21}. This date and the accompanying assessment are supplied dated curation, not independently verified here. That assessment labels the problem open and says the general invariant-order question remains open, citing the report's distinction between partial and bi-invariant orders and its stated lack of a nonexistence proof even for \(\mathrm{SL}(3,\Z)\).
\item The registry labels the statement \texttt{corrected\_verified} and says its mathematical content was checked against the official report, with notation, context, and question form normalized without changing the source-level claim. This is a supplied verification claim; no source visit or independent verification is asserted in this draft.
\end{itemize}
\end{document}
